\input{template.tex}

% Latin Modern has no T2A Cyrillic shapes; use Computer Modern for pdfLaTeX text.
\ifPDFTeX
  \renewcommand{\rmdefault}{cmr}
\fi

\def\SheetNo{2}

\begin{document}
\sheettitle{Множества, функции, отношения.}
\begin{center}\textbf{Краткие решения}\end{center}
Обозначение: $\mathbb N=\{1,2,\ldots\}$.
\begin{problems}

\item Пусть $A_1,A_2,\ldots,A_n$ "- конечные множества. Докажите, что
\[
\Bigl|\,\bigcup_{i=1}^{n}A_i\Bigr|
=\sum_{i}|A_i|-\sum_{i<j}|A_i\cap A_j|+\sum_{i<j<k}|A_i\cap A_j\cap A_k|-\ldots+(-1)^{n-1}|A_1\cap A_2\cap\ldots\cap A_n|,
\]
то есть
\[
\Bigl|\,\bigcup_{i=1}^{n}A_i\Bigr|
=\sum_{\varnothing\neq I\subseteq\{1,\ldots,n\}}(-1)^{|I|-1}\Bigl|\,\bigcap_{i\in I}A_i\Bigr|.
\]
\par\smallskip\noindent\textbf{Решение.} Пусть элемент $x$ принадлежит ровно $r\ge1$ из данных множеств. Тогда он входит ровно в $\binom rk$ пересечений по $k$ множеств, поэтому его суммарный вклад в правую часть равен
\[
\sum_{k=1}^{r}(-1)^{k-1}\binom rk
=1-\sum_{k=0}^{r}(-1)^k\binom rk
=1-(1-1)^r=1.
\]
Таким образом, каждый элемент объединения учитывается в правой части ровно один раз, а элементы вне объединения не учитываются. Следовательно, правая часть равна мощности объединения.\par

\item
Пусть $M$ "- множество всех трехэлементных подмножеств натурального ряда. Рассмотрим отображение $f\colon M \to \mathbb{N}$, которое каждому трехэлементному множеству ставит в соответствие его второй по величине элемент. Является ли это отображение а)~инъекцией б)~сюръекцией?
\par\smallskip\noindent\textbf{Решение.}
\textbf{а)} Нет: $\{1,2,3\}\ne\{1,2,4\}$, но $f(\{1,2,3\})=f(\{1,2,4\})=2$.
\textbf{б)} Нет: второй по величине элемент трёхэлементного подмножества $\mathbb N$ не может равняться $1$.\par

\item Постройте пример:
\begin{enumerate}[label=\alph*)]
  \item рефлексивного, симметричного, но не транзитивного отношения;
  \item рефлексивного, антисимметричного, но не транзитивного отношения;
  \item отношения, которое не симметрично и не антисимметрично;
  \item отношения, которое симметрично и антисимметрично.
\end{enumerate}
\par\smallskip\noindent\textbf{Решение.} Ну просто потерпеть, послушать студентов.\par

\item Рассмотрим на множестве $\mathbb{R}$ бинарное отношение $R$. Чему равно $R\circ R$, если $R(x,y)$ означает:
\begin{enumerate}[label=\alph*)]
  \item $y=x+1$;
  \par\smallskip\noindent\textbf{Решение.} Условие $(x,y)\in R\circ R$ означает существование $z$ такого, что $z=x+1$ и $y=z+1$. Следовательно, $y=x+2$. Обратно, при $y=x+2$ подходит $z=x+1$.\par
  \item $x+y=1$;
  \par\smallskip\noindent\textbf{Решение.} Требуется существование $z$ такого, что $x+z=1$ и $z+y=1$. Это равносильно условию $y=x$; промежуточный элемент $z=1-x$.\par
  \item $xy=1$;
  \par\smallskip\noindent\textbf{Решение.} Условия $xz=1$ и $zy=1$ равносильны $y=x\ne0$, причём $z=1/x$. Таким образом, $R\circ R=\{(x,x)\mid x\in\mathbb R\setminus\{0\}\}$.\par
  \item $xy>1$.
  \par\smallskip\noindent\textbf{Решение.} Ответ: $xy>0$. Если существуют $xz>1$ и $zy>1$, то $xyz^2>1$, откуда $xy>0$. Обратно, при $xy>0$ положим $z=1/x+1/y$. Тогда $xz=1+x/y>1$ и $zy=1+y/x>1$, поэтому $(x,y)\in R\circ R$.\par
\end{enumerate}

\item Является ли отношение $R$, такое что $(a,b)R(c,d)$, если $a+d=b+c$ на $\mathbb{N}\times\mathbb{N}$ отношением эквивалентности?
\par\smallskip\noindent\textbf{Решение.} Да. Условие $a+d=b+c$ равносильно равенству $a-b=c-d$. Рефлексивность, симметричность и транзитивность отношения следуют из соответствующих свойств равенства.\par

\item Композиция отношений $R$ и $S$ это отношение $T=RS$, где $xTy$, если найдется $z$, такой что $xRz$ и $zSy$. Пусть $R$ и $S$ "- транзитивные отношения на $A$. Будет ли транзитивной их композиция?
\par\smallskip\noindent\textbf{Решение.} Нет. На множестве $A=\{1,2,3,4\}$ рассмотрим отношения
\[
R=\{(1,2),(3,4)\},\qquad S=\{(2,3),(4,1)\}.
\]
Оба отношения транзитивны, поскольку в каждом из них отсутствуют цепочки из двух последовательных пар. Однако $RS=\{(1,3),(3,1)\}$: имеем $1\,(RS)\,3$ и $3\,(RS)\,1$, но $(1,1)\notin RS$. Следовательно, композиция не транзитивна.\par

\item Дано множество $X$, $|X|=n$. Сколько существует:
\begin{enumerate}[label=\cyralph*)]
\item бинарных отношений на множестве $X$?
\par\smallskip\noindent\textbf{Решение.} $2^{n^2}$. Отношение задаётся произвольной матрицей из нулей и единиц размера $n\times n$.\par
\item рефлексивных?
\par\smallskip\noindent\textbf{Решение.} $2^{n^2-n}$. На главной диагонали стоят единицы, остальные $n^2-n$ значений выбираются независимо.\par
\item симметричных?
\par\smallskip\noindent\textbf{Решение.} $2^{\,n+(n^2-n)/2}$. Независимо выбираются $n$ диагональных значений и по одному значению для каждой из $\binom n2$ пар симметричных позиций.\par
\item антисимметричных?
\par\smallskip\noindent\textbf{Решение.} $2^n\cdot3^{(n^2-n)/2}$. В каждой паре позиций $(i,j)$ и $(j,i)$ при $i\ne j$ допустимы ровно три варианта: $(1,0)$, $(0,1)$ и $(0,0)$. На диагонали значения произвольны.\par
\end{enumerate}

\item Какие из равенств верны для любых множеств $A,B,C$?
\begin{enumerate}[label=\cyralph*)]
  \item $(A\cap B)\cup C=(A\cup C)\cap(B\cup C)$;
  \par\smallskip\noindent\textbf{Решение.} Да.\par
  \item $(A\cup B)\cap C=(A\cap C)\cup(B\cap C)$;
  \par\smallskip\noindent\textbf{Решение.} Да.\par
  \item $(A\cup B)\setminus C=(A\setminus C)\cup B$;
  \par\smallskip\noindent\textbf{Решение.} Нет.\par
  \item $(A\cap B)\setminus C=(A\setminus C)\cap B$;
  \par\smallskip\noindent\textbf{Решение.} Да.\par
  \item $A\setminus(B\cup C)=(A\setminus B)\cap(A\setminus C)$;
  \par\smallskip\noindent\textbf{Решение.} Да.\par
  \item $A\setminus(B\cap C)=(A\setminus B)\cup(A\setminus C)$.
  \par\smallskip\noindent\textbf{Решение.} Да.\par
\end{enumerate}

\item Проведите подробное доказательство верных равенств предыдущей задачи, исходя из определений. (Докажем, что множества в левой и правой частях равны. Пусть $x$ "- любой элемент левой части равенства. Тогда\ldots{} Поэтому $x$ входит в правую часть. Обратно, пусть\ldots) Приведите контрпримеры к неверным равенствам.
\par\smallskip\noindent\textbf{Решение.} Ну просто потерпеть, послушать студентов. \par

\item На какое наименьшее число групп можно разбить все пары элементов данного $n$-элементного множества так, чтобы любые две пары из одной группы имели общий элемент?
\par\smallskip\noindent\textbf{Решение.} При $n\ge3$ ответ равен $n-2$; при $n=1$ ответ $0$, при $n=2$ - $1$.
\par Для построения занумеруем элементы числами $1,\ldots,n$. Для каждого $i=1,\ldots,n-3$ объединим в одну группу пары $\{i,j\}$ с $j>i$. Оставшиеся три пары на элементах $n-2,n-1,n$ образуют последнюю группу. Всего получаем $n-2$ группы.
\par Заметим, что любая группа попарно пересекающихся пар либо имеет общий элемент (\textit{звезда}), либо состоит из трёх пар $\{a,b\}$, $\{a,c\}$, $\{b,c\}$ (\textit{треугольник}). Действительно, если группа содержит $\{a,b\}$ и $\{a,c\}$, но не все её пары содержат $a$, то в ней есть $\{b,c\}$. Никакая другая пара не пересекает все три указанные пары. Группы из одной или двух пар имеют общий элемент.
\par Нижнюю оценку докажем индукцией по $n$. При $n=3$ требуется хотя бы одна группа. Пусть $n\ge4$ и имеется разбиение на $k$ групп. Если некоторая группа имеет общий элемент $x$, удалим $x$ и все пары, содержащие его. Эта группа исчезнет, поэтому пары оставшихся $n-1$ элементов разбиты не более чем на $k-1$ групп. По предположению индукции $k-1\ge n-3$, откуда $k\ge n-2$.
\par Если же ни одна группа не имеет общего элемента, все группы являются треугольниками и содержат по три пары. Поэтому
\[
k=\frac{\binom n2}{3}=\frac{n(n-1)}6\ge n-2,
\]
поскольку последнее неравенство равносильно $(n-3)(n-4)\ge0$.\par

\item Докажите, что в любой выборке из 52 положительных целых чисел найдутся хотя бы два, у которых либо их сумма, либо их разность делится на 100.
\par\smallskip\noindent\textbf{Решение.} Если два остатка по модулю $100$ совпадают, разность соответствующих чисел делится на $100$. Иначе все остатки различны. Разобьём их возможные значения на $51$ группу:
\[
\{0\},\quad\{50\},\quad\{r,100-r\}\quad(1\le r\le49).
\]
Среди $52$ различных остатков два попадут в одну группу; их сумма равна $100$.\par

\item Пусть $M$ "- количество всевозможных способов расставить слонов на шахматной доске так, чтобы никакие два из них не били друг друга. (Количество слонов может быть любым натуральным числом, даже 1.) Докажите, что $M+1$ "- квадрат некоторого натурального числа.
\par\smallskip\noindent\textbf{Решение.} Пусть $k$ - число допустимых расстановок на белых полях, включая пустую. Для чёрных полей число вариантов такое же по симметрии доски. Слоны на полях разных цветов не бьют друг друга, поэтому эти расстановки выбираются независимо. Исключая единственную пустую расстановку на всей доске, получаем $M=k^2-1$, откуда $M+1=k^2$.\par

\item У конечного множества $A$ выбрано несколько подмножеств. Каждое подмножество состоит хотя бы из двух элементов, а любые два пересекающихся подмножества имеют хотя бы два общих элемента. Докажите, что элементы множества $A$ можно покрасить в два цвета так, чтобы каждое из выбранных подмножеств содержало элементы обоих цветов.
\par\smallskip\noindent\textbf{Решение.} Индукция по числу выбранных подмножеств. Для одного подмножества окрасим два его элемента в разные цвета, а остальные элементы $A$ - произвольно. Пусть $A$ уже окрашено так, что первые $k$ подмножеств содержат оба цвета, и добавим подмножество $B$. Если в $B$ представлены оба цвета, ничего менять не требуется.
\par Иначе пусть все элементы $B$ красные. Перекрасим произвольный элемент $x\in B$ в синий. Поскольку $|B|\ge2$, теперь в $B$ есть оба цвета. Любое прежнее подмножество $C$, содержащее $x$, имеет с $B$ хотя бы два общих элемента, поэтому в нём остаётся красный элемент $y\in(B\cap C)\setminus\{x\}$. Синие элементы $C$ сохраняются. Подмножества, не содержащие $x$, не изменяются. Следовательно, все $k+1$ подмножеств содержат оба цвета.\par

\item Сколько различных выражений для множеств можно составить из переменных $A$ и $B$ с помощью (многократно используемых) операций пересечения, объединения и разности? (Два выражения считаются одинаковыми, если они равны при любых значениях переменных.) Тот же вопрос для трёх множеств и для $n$ множеств. (Ответ в общем случае: $2^{2^n-1}$.)
\par\smallskip\noindent\textbf{Решение.} Для фиксированного элемента $x$ обозначим через $a_i\in\{0,1\}$ признак его принадлежности множеству $A_i$. Операциям пересечения, объединения и разности соответствуют булевы операции $a\land b$, $a\lor b$ и $a\land\neg b$. Поэтому каждое выражение задаёт булеву функцию от $n$ переменных. Выражения равны при любых множествах тогда и только тогда, когда совпадают соответствующие функции: любой набор признаков принадлежности можно реализовать выбором множеств.
\par Все три операции сохраняют ноль, поэтому любая полученная функция удовлетворяет условию $f(0,\ldots,0)=0$. В таблице истинности значение на одной из $2^n$ строк фиксировано, а на остальных $2^n-1$ строках возможны два значения. Следовательно, различных выражений не более $2^{2^n-1}$.
\par Покажем, что каждая функция, сохраняющая ноль, реализуется. Для каждого непустого $I\subseteq\{1,\ldots,n\}$ выражение
\[
E_I=\left(\bigcap_{i\in I}A_i\right)\setminus\left(\bigcup_{j\notin I}A_j\right)
\]
выделяет элементы, принадлежащие ровно множествам с номерами из $I$. Если $I=\{1,\ldots,n\}$, вычитание опускаем. Объединяя $E_I$ для тех наборов, на которых функция равна единице, получаем требуемое выражение. Тождественный ноль задаётся выражением $A_1\setminus A_1$.
\par Таким образом, ответ равен $2^{2^n-1}$. Для двух множеств получаем $8$, для трёх - $128$.\par

\item Докажите, что $f^{-1}\bigl(g^{-1}(Z)\bigr)=(g\circ f)^{-1}(Z)$ для любого $Z\subseteq C$ и любых функций $f$ из множества $A$ в множество $B$ и $g$ из множества $B$ в множество $C$. Здесь $f^{-1}$ обозначает полный прообраз множества.
\par\smallskip\noindent\textbf{Решение.} Для любого $x\in A$ по определению прообраза
\[
\begin{aligned}
x\in f^{-1}\bigl(g^{-1}(Z)\bigr)
&\iff f(x)\in g^{-1}(Z)\\
&\iff g(f(x))\in Z\\
&\iff x\in(g\circ f)^{-1}(Z).
\end{aligned}
\]
Следовательно, множества равны.\par

\item На плоскости нарисовано некоторое количество равносторонних треугольников. Они не пересекаются, но могут иметь общие участки сторон. Мы хотим покрасить каждый треугольник в какой-нибудь цвет так, чтобы те из них, которые соприкасаются, были покрашены в разные цвета (треугольники, имеющие одну общую точку, могут быть покрашены в один цвет). Хватит ли для такой раскраски двух цветов?
\par\smallskip\noindent\textbf{Решение.} Да. Построим граф, вершины которого соответствуют треугольникам, а рёбра соединяют пары с общим участком стороны положительной длины. Ориентацию треугольника зададим направлением от его центра к любой вершине, взятым по модулю $120^\circ$; выбор вершины на это значение не влияет.
\par У соседних треугольников направления от центров к вершинам, противоположным общей стороне, взаимно противоположны. Поэтому при переходе по ребру ориентация меняется на $180^\circ\equiv60^\circ\pmod{120^\circ}$. В цикле длины $\ell$ получаем $60^\circ\ell\equiv0\pmod{120^\circ}$, откуда $\ell$ чётно. Граф без нечётных циклов двудолен, следовательно, двух цветов достаточно.\par

\item Из натуральных чисел от 1 до $2n$ выбрали $n+1$ число. Докажите, что одно из этих чисел делится на другое.
\par\smallskip\noindent\textbf{Решение.} Каждое выбранное число единственным образом представляется как $2^kq$, где $k\ge0$, а $q$ нечётно. Возможны лишь $n$ значений $q$: $1,3,\ldots,2n-1$. По принципу Дирихле у двух выбранных чисел нечётная часть одинакова. Тогда число с меньшим показателем степени двойки делит число с большим.\par

\item Верно ли для любых $f$, $Y_1$, $Y_2$ равенство $f^{-1}(Y_1\cap Y_2)=f^{-1}(Y_1)\cap f^{-1}(Y_2)$?
\par\smallskip\noindent\textbf{Решение.} Да. По определению прообраза
\[
\begin{aligned}
x\in f^{-1}(Y_1\cap Y_2)
&\iff f(x)\in Y_1\ \text{и}\ f(x)\in Y_2\\
&\iff x\in f^{-1}(Y_1)\cap f^{-1}(Y_2).
\end{aligned}
\]
Эта эквивалентность даёт оба включения.\par

\item Пусть $f$ всюду определена, а $g$ не всюду определена. Верно ли, что композиция $g\circ f$ не всюду определена?
\par\smallskip\noindent\textbf{Решение.} Нет. Пусть $f(x)=1$ для всех $x\in\mathbb R$, а $g(y)=1/y$ при $y\ne0$ и не определена при $y=0$. Тогда $(g\circ f)(x)=g(1)=1$ для любого $x\in\mathbb R$, то есть композиция всюду определена.\par

\item Докажите, что для любой функции $f\colon A\to A$, отображающей конечное множество $A$ в себя, найдутся различные целые положительные $m$, $n$, при которых $f^m$ и $f^n$ совпадают.
\par\smallskip\noindent\textbf{Решение.} При $|A|=N\ge1$ существует $N^N$ отображений $A\to A$. Все степени $f,f^2,\ldots$ - такие отображения, поэтому какие-то две положительные степени совпадают. При $A=\varnothing$ отображение единственно, поэтому любые две его степени совпадают.\par

\item Мы уже определили $f^n$ при целом $n\geqslant 2$, а также $f^{-1}$ для биекций. Как надо определить $f^0$, $f^1$ и отрицательные степени (для биекций), чтобы выполнялось равенство
\[
f^{m+n}=f^m\circ f^n
\]
(при любых целых $m$ и $n$)?
\par\smallskip\noindent\textbf{Решение.} Положим $f^0=\mathrm{id}_A$ и $f^1=f$. Для биекции $f\colon A\to A$ обозначим обратное отображение через $g=f^{-1}$ и при $k\ge1$ определим $f^{-k}=g^k$. Поскольку $f\circ g=g\circ f=\mathrm{id}_A$, имеем $f^k\circ g^k=\mathrm{id}_A$. Закон $f^{m+n}=f^m\circ f^n$ следует из ассоциативности композиции: для степеней одного знака складывается число множителей, а для степеней разных знаков взаимно обратные множители сокращаются.\par

\item Найдите количество:
\begin{enumerate}[label=\cyralph*)]
  \item неубывающих функций $f\colon\{1,2,\ldots,n\}\to\{1,2,\ldots,m\}$;
  \par\smallskip\noindent\textbf{Решение.} Выберем $n$ значений из $\{1,\ldots,m\}$ с повторениями. Расположив их в неубывающем порядке, однозначно получим последовательность $f(1),\ldots,f(n)$. Следовательно, число функций равно числу сочетаний с повторениями:
  \[
  \binom{m+n-1}{n}.
  \]\par
  \item неубывающих инъекций $f\colon\{1,2,\ldots,n\}\to\{1,2,\ldots,m\}$;
  \par\smallskip\noindent\textbf{Решение.} $\binom mn$ при $n\le m$; при $n>m$ ответ $0$. Нужно выбрать $n$ различных значений из $m$, после чего их возрастающий порядок единственным образом задаёт функцию. При $n\le m$
  \[
  \binom mn=\frac{m(m-1)\cdots(m-n+1)}{n!}.
  \]\par
  \item неубывающих сюръекций $f\colon\{1,2,\ldots,n\}\to\{1,2,\ldots,m\}$.
  \par\smallskip\noindent\textbf{Решение.} При $n<m$ сюръекций нет. Пусть $n\ge m$. Последовательность $f(1),\ldots,f(n)$ состоит из $m$ непустых последовательных блоков со значениями $1,2,\ldots,m$. Выберем $m-1$ границ блоков из $n-1$ промежутков между соседними позициями. Каждый выбор однозначно задаёт функцию, поэтому ответ равен
  \[
  \binom{n-1}{m-1}.
  \]\par
\end{enumerate}
Функция неубывающая, если $x\leqslant y$ влечёт $f(x)\leqslant f(y)$.

\item Докажите, что среди чисел Фибоначчи найдется число, оканчивающееся тремя нулями.
\par\smallskip\noindent\textbf{Решение.} Рассмотрим пары $(F_k,F_{k+1})$ по модулю $1000$, начиная с $(F_0,F_1)=(0,1)$. Возможных пар не более $1000^2$. Переход $(a,b)\mapsto(b,a+b)$ обратим: предыдущая пара восстанавливается как $(b-a,a)$. Поэтому из повторения какой-либо пары следует возврат к $(0,1)$ через некоторое число шагов $T>0$. Следовательно, $F_{jT}\equiv0\pmod{1000}$ при всех $j\ge1$. Выбрав $j$ достаточно большим, получим $F_{jT}>1000$.\par

\item В школе действует несколько кружков. Среди них есть кружок по топологии, который никто не посещает, и кружок по обществознанию, на который ходят все 2020 учеников школы. Списочные составы любых двух кружков различны. Кроме того, для любых двух кружков $A$ и $B$
\begin{enumerate}[label=\arabic*)]
  \item найдётся кружок $C$, который посещают в точности те ученики, которые посещают как $A$, так и $B$;
  \item найдётся кружок $D$, который посещают в точности те ученики, которые посещают хотя бы один из кружков $A$ или $B$.
\end{enumerate}
Докажите, что какой-то ученик посещает не менее половины всех кружков.
\par\smallskip\noindent\textbf{Решение.} Выберем непустой кружок $B$ с наименьшим числом участников. Для любого кружка $A$ пересечение $A\cap B$ также является кружком. По минимальности $B$ оно либо пусто, либо равно $B$. Таким образом, каждый кружок либо не пересекается с $B$, либо содержит $B$ целиком.
\par Каждому кружку $A$, не пересекающемуся с $B$, сопоставим кружок $A\cup B$. Он содержит $B$, причём разным $A$ соответствуют разные объединения, поскольку $A=(A\cup B)\setminus B$. Следовательно, кружков, содержащих $B$, не меньше, чем кружков, не пересекающихся с $B$. Каждый участник $B$ посещает все кружки первой группы, то есть не менее половины всех кружков.\par

\item Дано натуральное число $n$, и множество $S=\{1,2,\ldots,n\}$. Найдите наименьшее значение величины $|S\triangle A|+|S\triangle B|+|S\triangle C|$, где $A$ и $B$ "- непустые конечные множества натуральных чисел и $C=A+B=\{a+b \mid a\in A,\, b\in B\}$.
\par\smallskip\noindent\textbf{Решение.} Ответ: $n+1$.
\par Сначала докажем, что $|C|\ge |A|+|B|-1$. Запишем $A=\{a_1<\cdots<a_r\}$ и $B=\{b_1<\cdots<b_s\}$. Все $r+s-1$ элементов цепочки
\[
a_1+b_1<a_2+b_1<\cdots<a_r+b_1<a_r+b_2<\cdots<a_r+b_s
\]
различны и принадлежат $C$, что даёт требуемую оценку (или можно просто по индукции по размеру А).
\par Обозначим искомую сумму через $D$. Раскрывая мощности симметрических разностей, получаем
\[
\begin{aligned}
D&=3n+|A|+|B|+|C|-2|S\cap A|-2|S\cap B|-2|S\cap C|\\
 &\ge 3n+2|A|+2|B|-1-2|S\cap A|-2|S\cap B|-2|S\cap C|\\
 &\ge 3n-1-2|S\cap C|.
\end{aligned}
\]
Здесь использованы неравенства $|S\cap A|\le |A|$ и $|S\cap B|\le |B|$. Поскольку элементы $A$ и $B$ положительны, $1\notin C$. Следовательно, $|S\cap C|\le n-1$, откуда
\[
D\ge3n-1-2(n-1)=n+1.
\]
Равенство достигается при $A=\{1\}$ и $B=S$: тогда $C=\{2,3,\ldots,n+1\}$ и $D=(n-1)+0+2=n+1$.\par

\end{problems}
\end{document}
